% SPDX-License-Identifier: LPPL-1.3c
% Copyright (C) 2026 Christophe Jorssen
% Author and Current Maintainer: Christophe Jorssen <christophe.jorssen@gmail.com>
% LPPL maintenance status: maintained
% This file is part of luafemm. See LICENSE and LICENSES.md for its terms.
\section{A small excitation around a saturated operating point}
\label{sec:operating}

\subsection{Secant and differential inductance}
The flux-linkage curve bends as the iron saturates. The inductance seen
by a small additional excitation depends on the local slope of that curve.
Dividing flux linkage by current gives a secant slope; a small perturbation
sees the tangent slope. These are different questions, even when the geometry
and the winding have not changed.

An \emph{operating point} captures a converged static solution and its exact
ordered mesh. A second picture inherits the geometry and material declarations,
but uses a linear constitutive law evaluated at that saved state. Two modes
are available: |frozen| uses the local secant permeability, while
|incremental| uses the differential tensor. Both planar and axisymmetric
static models are supported. Section~\ref{sec:biased-harmonic} extends incremental
mode to a small harmonic excitation. Neither mode advances time or moves the rotor.

\subsection{First solve, then capture, then perturb}
Declare materials, boundaries and windings in an ordinary |femm/problem|,
then schedule capture with a picture-level key:
\begin{codeexample}[code only]
\tikzset{biased coil/.style={
  femm/problem={xmin=-30,xmax=30,ymin=-25,ymax=25,depth=20},
  femm/circuit={name=drive,current=5},
  femm/operating point={name=bias},
  femm/cache={file=coil-bias,mode=auto}}}
\end{codeexample}
The |femm/operating point| key schedules the capture at the end of that
picture, after all its physical paths have been declared. It solves if
necessary, then takes an independent copy. The snapshot is available only
after the picture has ended. Alternatively,
|\femmoperatingpoint[name=bias]| performs the capture immediately at its
position in the input; place it after the last physical declaration.
A capture does not replace the current model.

The next picture selects the saved name. It does \emph{not} repeat the
physical paths, material definitions or circuit declarations:
\begin{codeexample}[code only]
\begin{tikzpicture}[
  femm/problem={operating point=bias,permeability=incremental,
    circuit currents={drive=1}},
  femm/cache={file=coil-increment,mode=auto}]
  \pic {femm geometry};
  \pic[femm/field style={orange!75!black}] {femm field};
  \femmmeasure[name=small,quantity=circuit,circuit=drive]
\end{tikzpicture}
\end{codeexample}
Here the current is $\delta I=1\,\mathrm A$. This is a normalisation of a
linear response: the resulting linkage is numerically the differential
inductance in henries. It is not a claim that the nonlinear circuit stays
linear over a finite one-ampere excursion. For an actual small change $\epsilon$,
$\mathbf B(I_0+\epsilon)\simeq\mathbf B_0+
\epsilon\,\partial\mathbf B/\partial I$.

In incremental mode every independent bias source is zeroed: circuit currents,
impressed current densities, coercive sources and imposed boundary loads.
The specified |circuit currents| then supply the perturbation. Unlisted
circuits have zero perturbation current. Robin coefficients, boundary kinds
and periodic ties are inherited; their additive loads are homogeneous.
The resulting lines and profiles describe $\delta\mathbf B$ and
$\delta\mathbf H$, not the total field. Ordinary |femm/profile|,
|\femmplot|, |\femmmeasure| and |\femmresultplot| keep their usual syntax.

To reconstruct the bias, select frozen permeability:
\begin{codeexample}[code only]
\begin{tikzpicture}[
  femm/problem={operating point=bias,permeability=frozen}]
  \pic {femm geometry};
  \pic {femm field};
\end{tikzpicture}
\end{codeexample}
Frozen mode retains every bias source by default. With unchanged sources its
solution recovers the nonlinear bias to solver tolerance. An explicit circuit
current replaces that circuit's current; it is not added to it. This permits
source decomposition when all other sources are accounted for. Retained
magnets or nonzero boundary loads must not be mistaken for zero excitation.

\subsection{One core, two slopes}
The following worked example compares the bias and incremental fields of a
closed rectangular core, then plots its flux linkage and both inductances.
The B--H table is illustrative, not a measured steel specification. The field
panels deliberately use separate contour levels; counting lines across the
two panels does not measure their flux ratio.

Start with |\Mesh=3| while arranging the figures, then try |1.5| and a finer
value before interpreting the curves quantitatively. Each bias solve has its
own mesh, but its frozen and incremental responses reuse that mesh exactly.
Within a bias-response pair there is no remeshing or transfer error. Across a
current sweep, geometry-driven meshing is deterministic for unchanged inputs.
Convergence of the tangent solve on one mesh does not establish convergence
of the spatial discretisation.

\luafemmexamplebegin{operating-point}

The common geometry and illustrative B--H law define every bias. The sweep captures each solution before requesting its differential response.

% manual example: operating-point/1

\begin{codeexample}[setup code]
\definecolor{fluxblue}{RGB}{20,85,135}
\definecolor{response}{RGB}{185,65,45}
\newbox\BiasCalculation
\def\Mesh{1.5}
\tikzset{
  femm/materials/core/.style={
    bh={0/0,.25/25.293,.5/59.375,.75/146.191,1/400,
      1.25/1040.527,1.5/2428.125,1.75/5098.926,2/9800,
      2.25/17524.512,2.5/29546.875,3/73200}},
  bias model/.style={femm/problem={xmin=-30,xmax=30,ymin=-25,ymax=25,
    mesh size=\Mesh,depth=20,interpolation=femm},
    femm/circuit={name=drive,current=#1}},
  coil leg/.style={fill=orange!30,draw=orange!65!black,
    femm/region={material=copper,current density=0,
      circuit=drive,turns=#1,winding fill=.65}},
  response model/.style={femm/problem={operating point=bias-#1,
    permeability=incremental,circuit currents={drive=1}}},
  femm/arrow spacing=12mm}
\newcommand\CoreGeometry{
  \filldraw[even odd rule,fill=blue!7,draw=fluxblue,
    femm/region={material=core}]
    (-15,-15) rectangle (15,15) (-8,-8) rectangle (8,8);
  \filldraw[coil leg=100] (-21,-7) rectangle (-17,7);
  \filldraw[coil leg=-100] (-6,-7) rectangle (-2,7);
  \node at (-19,0) {$\odot$}; \node at (-4,0) {$\otimes$};}
\end{codeexample}

The following definitions complete the shared setup.

% manual example: operating-point/2

\begin{codeexample}[setup code]
\foreach \Index/\Bias in {a/.25,b/.5,c/1,d/2,e/3,f/5,g/8,h/12}{
  \setbox\BiasCalculation=\hbox{
    \begin{tikzpicture}[bias model=\Bias,
      femm/cache={file=\jobname-bias-\Index,mode=auto},
      femm/operating point={name=bias-\Index}]
      \CoreGeometry
      \femmmeasure[name=dc-\Index,quantity=circuit,circuit=drive,parameter=\Bias]
    \end{tikzpicture}}
  \setbox\BiasCalculation=\hbox{
    \begin{tikzpicture}[response model=\Index,
      femm/cache={file=\jobname-response-\Index,mode=auto}]
      \femmmeasure[name=small-\Index,quantity=circuit,circuit=drive,parameter=\Bias]
    \end{tikzpicture}}}
\end{codeexample}

The following picture uses the preceding definitions.

% manual example: operating-point/3

\begin{codeexample}[width=0pt]
\begin{tikzpicture}[femm/problem={operating point=bias-f,permeability=frozen},
  femm/cache={file=\jobname-frozen,mode=auto},scale=1.05]
  \clip (-25,-20) rectangle (22,23);
  \pic[femm/geometry style={draw=fluxblue!60}] {femm geometry};
  \pic[femm/lines=19,femm/field style={fluxblue}] {femm field};
  \node[fill=white,font=\small,align=center] at (0,20)
    {Bias field\\$I_0=5\,\mathrm A$};
  \node[font=\scriptsize,fill=white] at (0,-18) {Frozen reconstruction};
\end{tikzpicture}
\end{codeexample}

% manual example: operating-point/4

\begin{codeexample}[width=0pt]
\begin{tikzpicture}[response model=f,
  femm/cache={file=\jobname-response-f,mode=auto},scale=1.05]
  \clip (-25,-20) rectangle (22,23);
  \pic[femm/geometry style={draw=black!45}] {femm geometry};
  \pic[femm/lines=19,femm/field style={response}] {femm field};
  \node[fill=white,font=\small,align=center] at (0,20)
    {Incremental response\\$\delta I=1\,\mathrm A$};
  \node[font=\scriptsize,fill=white] at (0,-18) {Tangent field, not total field};
\end{tikzpicture}
\end{codeexample}

% manual example: operating-point/5

\begin{codeexample}[setup code]
\def\BiasResults{dc-a,dc-b,dc-c,dc-d,dc-e,dc-f,dc-g,dc-h}
\def\SmallResults{small-a,small-b,small-c,small-d,small-e,small-f,small-g,small-h}
\end{codeexample}

The following picture uses the preceding definitions.

% manual example: operating-point/6

\begin{codeexample}[width=0pt]
\begin{tikzpicture}
\begin{axis}[width=7.5cm,height=5.6cm,grid=major,
  xlabel={Bias current $I_0$ (A)},ylabel={Flux linkage $\lambda_0$ (mWb)},
  y filter/.expression={1000*y},scaled y ticks=false,
  title={The core saturates}]
  \femmresultplot[results={\BiasResults},component=linkage,
    plot options={fluxblue,thick,mark=*}]
\end{axis}
\end{tikzpicture}
\end{codeexample}

% manual example: operating-point/7

\begin{codeexample}[width=0pt]
\begin{tikzpicture}
\begin{axis}[width=7.5cm,height=5.6cm,grid=major,
  xlabel={Bias current $I_0$ (A)},ylabel={Inductance (mH)},
  y filter/.expression={1000*y},scaled y ticks=false,
  title={Two different inductances},
  legend style={font=\scriptsize,at={(.98,.98)},anchor=north east}]
  \femmresultplot[results={\BiasResults},component=inductance,
    plot options={fluxblue,thick,mark=*}]
  \addlegendentry{$\lambda_0/I_0$ (secant)}
  \femmresultplot[results={\SmallResults},component=incremental_inductance,
    plot options={response,thick,mark=square*}]
  \addlegendentry{$\mathrm d\lambda/\mathrm dI$ (incremental)}
\end{axis}
\end{tikzpicture}
\end{codeexample}

\luafemmexampleend{operating-point}

The result component |inductance| is $\lambda/I$ for the ordinary static
model under its existing eligibility rules. On an incremental response with
one nonzero circuit excitation it is $\delta\lambda/\delta I$; the explicit
alias |incremental_inductance| names the same value. Use the alias in
plots to make the interpretation unambiguous. Its SI unit is the henry:
\begin{codeexample}[code only]
\femmvalue[component=incremental_inductance]{small}
\femmresultplot[results={small-a,small-b,small-c},
  component=incremental_inductance,abscissa=parameter,
  plot options={orange!75!black,thick,mark=square*}]
\end{codeexample}
Record |parameter=| as the bias current, as the complete example does.
The result's |current| is the perturbation current, not the bias current.

For several windings, excite one port at a time and measure the linkage of
all ports. With a unit excitation at port $j$, the linkage at port $i$ gives
$L_{ij}$. Unexcited ports still report linkage but cannot report a self
inductance by division by zero. When several circuits are excited together,
the response does not report an |inductance| ratio: use the individual
linkages and the declared excitation vector. A matrix is assembled by the
caller; there is no automatic matrix-valued result or circuit-network solver.
For frozen responses the ratio is also omitted in the presence of other
independent sources. In either mode a planar winding must be closed in signed
turns for its linkage to be gauge independent.

\subsection{Key reference}
\begin{key}{/tikz/femm/operating point=\marg{options}}
Schedule a capture at the end of the originating picture. Options belong to
|/femm/operating point|. Use once for each named snapshot; names are
unique across the document. Reusing a name is an error, which prevents a
later picture from silently changing the meaning of an earlier response.
\end{key}
\begin{command}{\femmoperatingpoint\oarg{options}}
Capture immediately, solving first if required. This command and the picture
key have the same options. A capture owns its data and does not borrow a
mutable potential array from the original model.
\end{command}
\begin{key}{/femm/operating point/name=\meta{name} (initially empty)}
Required nonempty snapshot identifier. Prefer letters, digits and hyphens in
examples. A snapshot name is separate from a measurement or profile name.
\end{key}
\begin{key}{/femm/operating point/model=\meta{model name} (initially empty)}
Select a registered |/femm/problem/name| for immediate capture; empty means
the current model. The picture-level key supplies its own originating model
identifier automatically, so another picture cannot redirect that capture.
\end{key}
\begin{key}{/femm/problem/operating point=\meta{name} (initially empty)}
Start a response from a previously captured operating point. Empty starts an
ordinary problem. Geometry, units, problem type, material laws, solver settings,
depth and mesh all come from the snapshot. Explicit geometry and outer-solver
options in the same problem declaration are rejected. Set those on the bias
problem. The six linear-solver keys in Section~\ref{sec:linear-controls} may be
explicitly overridden on the response; otherwise their captured values persist. |name| may name the new response model;
|femm/cache| and presentation keys remain independent. The one physical override
is |frequency|, for the incremental harmonic extension described next.
\end{key}
\begin{key}{/femm/problem/permeability=\meta{choice} (initially incremental)}
A PGF choice, either |incremental| or |frozen|. It selects the linear material
law and the source defaults explained above. It is distinct from
|/femm/cache/mode=frozen|, which only requires an existing disk cache.
\end{key}
\begin{key}{/femm/problem/circuit currents=\marg{name=value list} (initially empty)}
Replace selected circuit currents, in amperes, in the inherited problem.
For example |circuit currents={drive=1,sense=0}|. PGF evaluates the numerical
expressions, so |drive=1/2| is allowed. Circuit names must already exist in
the bias. An empty list retains the mode's defaults; in incremental mode
this gives a zero response. Use this key instead of redeclaring inherited
circuits with |femm/circuit|. The enclosing problem remains styleable:
\begin{codeexample}[code only]
\tikzset{unit response/.style={femm/problem={
  operating point=#1,permeability=incremental,
  circuit currents={drive=1}}}}
\end{codeexample}
\end{key}

\subsection{What is actually linearised?}
Let the isotropic material table describe the noncoercive law
$\mathbf H_{\mathrm{tab}}=\nu(B)\mathbf B$, where $\nu(B)=H(B)/B$ in SI.
At the static bias $\mathbf B_0$, define $B_0=|\mathbf B_0|$ and
$\mathbf n=\mathbf B_0/B_0$. Frozen permeability uses the scalar reluctivity
$\nu_0=H(B_0)/B_0$. The incremental law is
\[
 \delta\mathbf H=\boldsymbol\nu_{\mathrm{inc}}\,\delta\mathbf B,
 \qquad
 \boldsymbol\nu_{\mathrm{inc}}=
 \nu_0\mathbf I+
 \left(\left.\frac{\mathrm dH}{\mathrm dB}\right|_{B_0}-\nu_0\right)
 \mathbf n\mathbf n^{\mathsf T}.
\]
Thus parallel and perpendicular perturbations see different slopes. At zero
bias the initial table slope is used in both directions. An already linear
tensor retains its tensor in both modes. Permanent-magnet coercivity belongs
to the source term, not to this derivative; incremental mode removes that
source while preserving the recoil law.

The implementation uses the same B--H interpolation and material tangent as
the nonlinear solver. At a corner of a piecewise-linear B--H table the
classical derivative is not unique: the selected interval's tangent is used.
Avoid interpreting a centred finite difference crossing such a corner as an
exact derivative check. Smooth FEMM-style interpolation is available, but
its smoothing changes the effective input curve as described in the material
chapter.

For planar P1 elements the bias curl and tangent are constant in each
triangle. For axisymmetry, $B_r=-\partial_zA_\phi$ and
$B_z=\partial_rA_\phi+A_\phi/r$ vary within an element. The bias and material
tangent are therefore evaluated at each volume quadrature point, with the
same radial weights as the nonlinear solve. Pointwise H profiles use the
local tensor at the sampled point. The static weak operator and its signed
boundary constraints are otherwise unchanged.

Regional |energy| and |coenergy| on an incremental response both mean the
quadratic response energy $\tfrac12\int\delta\mathbf B\cdot
\delta\mathbf H\,\mathrm dV$. They exclude the bias energy and the first-order
cross term in a total-energy expansion. For a single unit-current response,
$2W=L_{\mathrm{inc}}$. In frozen mode the energy refers to the frozen linear
material and retained coercive source; it is not the original nonlinear
stored energy even when the field reconstruction agrees.

\subsection{Persistence, checks and current limits}
Named operating points live for the TeX run. Across compilations, keep the
bias declarations and use their ordinary native solution cache: capture can
then copy the restored solution without repeating Newton iterations. Give
each response its own cache filename. Passive cache format 25 includes the
ordered bias mesh, nodal bias field, constitutive mode, source provenance and
response excitations. A changed bias cannot reuse an old response silently.
Older caches must be recomputed. There is no transfer between unrelated meshes
and supported static FEMM |.ans| snapshots can supply an owned bias; see
Section~\ref{sec:interchange}.

The operating-point model accepts native planar/axisymmetric static models with
single-valued isotropic nonlinear, directional DC sheet or positive linear tensor laws, including
supported permanent magnets. Named circuits must use prescribed series
currents. Source/boundary callbacks, ideal confined
models, voltage/parallel circuit drives and capture of an already linearised
response are rejected. Geometry and physics are locked after creating a
response: draw labels or untagged decorative paths, but do not add regions,
change materials or redeclare boundaries. Create a new nonlinear bias for a
changed geometry. Small harmonic responses are described in
Section~\ref{sec:biased-harmonic}. Mechanical first variations are described in
Section~\ref{sec:mechanical-response}. Planar FEM/ANS response exchange is described in
Section~\ref{sec:response-exchange}. Independent rotor positions are described in
Section~\ref{sec:sliding}.
The nonlinear bias itself retains its ordinary supported export capabilities.

The regression suite checks frozen reconstruction, the analytic oblique
material tensor, permanent-magnet source removal, zero/linear limits,
reciprocity, response energy and finite differences of planar and
axisymmetric nonlinear solves. The optional xfemm study compares against
independent ordinary nonlinear solves; the FEMM 5 previous-solution comparison additionally checks the bounded
planar response exchange. Detailed tolerances
and measured errors are recorded in \nolinkurl{doc/validation.md}.

\subsection{Lua entry points}
The numerical module |require("luafemm-operating")| does not access TeX globals.
|capture(model)| returns an owned snapshot of a converged static problem;
|new(snapshot,options)| constructs an unsolved response. Options are |mode|,
|currents| (a table from circuit names to amperes), |frequency| (hertz),
|phases| (a table from circuit names to degrees) and |cache| (the ordinary
native cache options). They default to |incremental|, empty current/phase tables,
zero frequency and caching off. Unknown options are rejected. For example:
\begin{codeexample}[code only]
local O = require("luafemm-operating")
local bias = O.capture(model) -- model has already converged
local response = O.new(bias, {
  mode = "incremental", currents = {drive = 1},
  cache = {file = "unit-response", mode = "auto"}
})
require("luafemm").solve(response)
local r = require("luafemm-static").result(response, "drive")
-- r.incremental_inductance is in henries when available.
\end{codeexample}
A snapshot may also come from a compatible restored native cache. The snapshot
itself is an in-memory object; |capture| does not create an extra disk format.
